\clearpage
\section{Supplementary Appendix S1: Prespecified review protocol}
\label{app:protocol}

This editable appendix summarizes the methods specified in the prospectively registered PROSPERO protocol for \textit{Conceptual Clarification of Digital Patient Frameworks, Methodological Approaches, and Infrastructural Gaps: Evidence from a Systematic Meta-Review}. The review was scheduled to begin on 24 October 2025. Preliminary searches and piloting of study selection had been completed at registration; formal screening had started, whereas extraction, appraisal, and synthesis had not.

\subsection{Aim and review questions}

The primary aim was to synthesize how the digital-patient spectrum---digital twins, virtual patients, in silico patients, and synthetic patients---is conceptualized, constructed, validated, and operationalized in healthcare, and to identify methodological and infrastructural gaps that limit clinical translation. The six prespecified questions addressed:

\begin{enumerate}[label=\textbf{RQ\arabic*.},leftmargin=*]
    \item definitional boundaries, distinguishing features, classifications, and hierarchy levels across the digital-patient spectrum;
    \item model architectures, data pipelines, coupling, temporal orientation, computational approaches, infrastructure, and interoperability standards;
    \item validation methods, fidelity measures, reproducibility, generalizability, and trustworthiness;
    \item healthcare domains, functional roles, maturity, personalization, benefits, and limitations;
    \item ethical, legal, technical, data, interoperability, organizational, and infrastructural barriers; and
    \item research directions, validation and reporting standards, emerging technologies, and governance priorities needed for clinical readiness.
\end{enumerate}

\subsection{Information sources and search concepts}

The protocol specified Web of Science, PubMed, IEEE Xplore, and Google Scholar, supplemented by reference-list searching and searches of TechRxiv, arXiv, bioRxiv, engrXiv, and medRxiv. Database syntax was to be adapted while preserving three concept blocks:

\begin{enumerate}[leftmargin=*]
    \item \textbf{Digital-patient construct:} digital twin, patient digital twin, digital patient, virtual patient, in silico patient, synthetic patient, virtual sibling, patient-specific computational model, virtual physiological human, VPH, or surrogate model.
    \item \textbf{Health context:} patient, medicine, healthcare, surgery, treatment, pharmacology, oncology, clinical care, physiology, pathophysiology, biomechanics, imaging, genomics, precision medicine, personalized medicine, medical devices, quantitative medicine, or preventive medicine.
    \item \textbf{Review design:} systematic review, scoping review, meta-analysis, meta-review, review of reviews, umbrella review, or structured review.
\end{enumerate}

Searches were limited to English-language reports published from 2020 onward. The protocol scheduled the initial search for October 2025 and an update before final manuscript submission. Preprints were to be checked for subsequently published versions to avoid duplicate inclusion.

\subsection{Eligibility criteria}

Eligible reports were systematic, scoping, umbrella, meta-, or structured narrative reviews that explicitly examined at least one generic or patient-specific digital model applied to human health. Models could represent biological organization from subcellular, cellular, tissue, and organ levels through whole-body and in silico populations. Eligible reviews had to address conceptual framing, model construction, application, validation, methodological evaluation, or scalability.

Exclusions comprised primary studies without a structured review methodology; editorials, perspectives, commentaries, and abstracts without sufficient methods; non-English reports; animal-only, non-biological, or purely technical models without patient linkage; and digital twins limited to hospitals, logistics, infrastructure, communication systems, or devices without patient-level modeling. Behavioral or educational simulations without a physiological, biomechanical, or biological patient representation were outside scope.

\subsection{Selection, extraction, and appraisal}

The protocol specified independent title and abstract assessment by two reviewers in CADIMA, followed by adjudication under the procedure in Supplementary Appendix S2. A standardized, piloted extraction form was organized around the six review questions. One reviewer extracted each included report and a second reviewer checked the entry for accuracy.

The protocol prespecified review-type-appropriate methodological appraisal using AMSTAR~2 or SANRA and a domain-based CADIMA assessment of review-level bias and validity. Appraisal covered eligibility, search completeness, study selection, extraction reliability, appraisal of included evidence, synthesis, validation claims, generalizability, implementation, and sponsorship. Ratings were intended to contextualize the synthesis rather than determine eligibility automatically.

\subsection{Synthesis}

The planned synthesis was narrative and descriptive, supported by tabular summaries and conceptual maps. Prespecified comparisons included generic or archetype-based versus individual patient-specific models; static versus dynamic or continuously coupled models; and educational applications versus clinical decision support and personalized care. No quantitative pooling was anticipated.

\clearpage
\section{Supplementary Appendix S2: Screening inconsistency-resolution protocol}
\label{app:adjudication}

\subsection{Title and abstract screening}

Each record was evaluated independently by two reviewers using three core questions:

\begin{enumerate}[leftmargin=*]
    \item Is the report a review with structured methodology?
    \item Does it address a digital patient, digital twin, or an eligible adjacent construct?
    \item Is it focused on human health or healthcare?
\end{enumerate}

Permitted responses were Yes, Unclear, and No. A record was considered potentially eligible for full-text review when there was no agreed No on any core question. Uncertainty at this stage therefore favored retrieval of the full report.

\subsection{Classification and resolution of discrepancies}

CADIMA-flagged inconsistencies were divided according to whether they could change the inclusion decision.

\paragraph{Category A: no effect on eligibility.}
This category included records for which both reviewers recorded No on at least one criterion, reviewers identified different but independently sufficient exclusion criteria, or a Yes/Unclear discrepancy occurred while another criterion already had two No judgments. The coordinator verified the pattern and finalized the exclusion as the third reviewer.

\paragraph{Category B: possible effect on eligibility.}
This category included direct Yes/No conflicts and mixed response patterns that could determine whether the report advanced. The coordinator initiated record-specific discussion; both reviewers briefly restated their rationales and indicated inclusion or exclusion. If uncertainty remained, the report advanced. When consensus could not be reached, the coordinator made the final decision as the third reviewer.

\subsection{Documentation responsibilities}

The coordinator categorized each discrepancy, facilitated Category B discussions, acted as the third reviewer when required, and ensured that the final judgment and rationale were documented in CADIMA. Reviewers provided concise, evidence-based rationales and participated in discussion for eligibility-affecting conflicts. This procedure was intended to make adjudication consistent, auditable, and conservative at the screening stage.
